% !TeX root = booklet-02.tex
\subsection{\lr{Mobility Degradation} و \lr{Velocity Saturation}}
\label{sec:high-field}
\subsubsection*{دو میدان با دو نقش فیزیکی}
میدان عمودی، حامل‌های متحرک را نزدیک فصل مشترک \lr{Silicon} و اکسید
نگه می‌دارد. افزایش این میدان، پراکندگی حامل‌ها در نزدیکی فصل مشترک را
شدیدتر می‌کند و \lr{Mobility} مؤثر کاهش می‌یابد. میدان طولی در امتداد کانال،
حرکت رانشی ایجاد می‌کند. در میدان طولی بزرگ، سرعت با افزایش میدان به یک
حد نزدیک می‌شود. این دو پدیده در کتاب و اسلایدها به‌ترتیب
\lr{Mobility Degradation} و \lr{Velocity Saturation} نامیده شده‌اند
\cite[صفحات چاپی ۷۴ تا ۷۶]{weste}\cite[اسلایدهای ۶ تا ۹]{nonideal}.

برای برآورد مرتبه بزرگی، اسلاید ۶ از دو رابطه استفاده می‌کند
\begin{equation}
 E_{\mathrm{vert}}\sim\frac{V_{GS}}{t_{\mathrm{ox}}},\qquad
 \overline E_{\mathrm{lat}}=\frac{V_{DS}}{L}.
 \label{eq:fields-estimate}
\end{equation}
رابطه اول نمایش تقریبی وابستگی میدان عمودی به بایاس و ضخامت است.
افت واقعی ولتاژ اکسید و میدان مؤثر حامل‌ها به ساختار و توزیع بار وابسته‌اند.
رابطه دوم، میانگین میدان طولی روی طول کانال را می‌دهد. در استخراج رابطه
جریان، میدان محلی \(E(x)\) را دنبال می‌کنیم.

\subsubsection*{مدل تجربی کاهش \lr{Mobility}}
در رابطه \secnum{۲.۲۳} کتاب و اسلاید ۸، \lr{Mobility} مؤثر با یک مدل تجربی
برحسب بایاس گیت و ضخامت اکسید نوشته شده است. برای \lr{nMOS} و
\lr{pMOS}، روابط را با قرارداد بزرگی مثبت ولتاژها چنین بازنویسی می‌کنیم
\begin{align}
 \mu_{\mathrm{eff},n}
 &=\frac{\qty{540}{\centi\metre\squared\per\volt\per\second}}
 {1+\left[\dfrac{V_{GS}+V_{tn}}
 {\bigl(\qty{0.54}{\volt\per\nano\metre}\bigr)t_{\mathrm{ox}}}\right]^{1.85}},
 \label{eq:mobility-n}\\[3pt]
 \mu_{\mathrm{eff},p}
 &=\frac{\qty{185}{\centi\metre\squared\per\volt\per\second}}
 {1+\dfrac{\left|V_{SG}+1.5V_{TP}\right|}
 {\bigl(\qty{0.338}{\volt\per\nano\metre}\bigr)t_{\mathrm{ox}}}}.
 \label{eq:mobility-p}
\end{align}
در \eqref{eq:mobility-p}، \(V_{TP}=|V_{tp}|\) است. کتاب در این رابطه،
ولتاژهای \lr{pMOS} را به‌صورت بزرگی مثبت به کار می‌برد
\cite[رابطه \secnum{۲.۲۳} و مثال \secnum{۲.۳}]{weste}. صورت هر کسر،
مقیاس \lr{Mobility} در همان مدل تجربی است. اعداد ۵۴۰ و ۱۸۵ به همراه
واحدشان پارامترهای مدل‌اند. این ضرایب از برازش داده‌های افزاره به دست آمده‌اند.

در مخرج \eqref{eq:mobility-n}، ضریب
\(\qty{0.54}{\volt\per\nano\metre}\) در ضخامت ضرب می‌شود و نتیجه ولتاژ
است. نسبت دو ولتاژ بی‌بعد است و توان \(1.85\) به همین نسبت اعمال می‌شود.
اگر ضخامت را با واحد \(\unit{\nano\metre}\) وارد می‌کنیم، مقدار عددی
\(1.05\) برای ضخامت \(\qty{1.05}{\nano\metre}\) به کار می‌رود.
برای محاسبه کاملاً \lr{SI}، ضریب متناظر
\(\qty{5.4e8}{\volt\per\metre}\) است.

با ثابت بودن آستانه و ضخامت، افزایش \(V_{GS}\) در دامنه هدایت
\lr{nMOS}، مخرج \eqref{eq:mobility-n} را بزرگ‌تر و \(\mu_{\mathrm{eff},n}\)
را کوچک‌تر می‌کند. این نتیجه، کاهش \lr{Mobility} در مدل را نشان می‌دهد.
جریان هم‌زمان به مقدار بار کانال وابسته است و افزایش گیت می‌تواند با وجود
کاهش \lr{Mobility}، جریان را افزایش دهد. سهم هر عامل باید در رابطه کامل
جریان بررسی شود.

\subsubsection*{حل مرحله‌به‌مرحله مثال \secnum{۲.۳} کتاب}
مثال کتاب از \(V_{GS}=V_{SG}=\qty{1.0}{\volt}\)،
\(V_{tn}=V_{TP}=\qty{0.3}{\volt}\) و
\(t_{\mathrm{ox}}=\qty{1.05}{\nano\metre}\) استفاده می‌کند
\cite[صفحه چاپی ۷۶]{weste}. ابتدا بخش بی‌بعد مخرج را حساب می‌کنیم
\begin{align}
 a_n&=\frac{1.0+0.3}{0.54\times1.05}\approx2.293,
 &a_n^{1.85}&\approx4.642,\notag\\
 a_p&=\frac{1.0+1.5(0.3)}{0.338\times1.05}\approx4.086.
 \label{eq:mobility-example-ratios}
\end{align}
جایگذاری این اعداد در دو رابطه \lr{Mobility} نتیجه می‌دهد
\begin{align}
 \mu_{\mathrm{eff},n}
 &\approx\frac{540}{1+4.642}
 \approx\qty{95.7}{\centi\metre\squared\per\volt\per\second},\notag\\
 \mu_{\mathrm{eff},p}
 &\approx\frac{185}{1+4.086}
 \approx\qty{36.4}{\centi\metre\squared\per\volt\per\second}.
 \label{eq:mobility-example-results}
\end{align}
اعداد گرد‌شده کتاب، به‌ترتیب ۹۶ و ۳۶ با واحد
\(\unit{\centi\metre\squared\per\volt\per\second}\) هستند. در مثال
ولتاژ بحرانی از همین گردکردن استفاده می‌کنیم. برای محاسبات چندمرحله‌ای،
نگهداری رقم‌های بیشتر تا انتهای محاسبه خطای گردکردن را محدود می‌کند.

مثال \secnum{۲.۱} کتاب برای مدل ایده‌آل یک مقدار ثابت
\(\qty{80}{\centi\metre\squared\per\volt\per\second}\) به \lr{Mobility}
الکترون نسبت می‌دهد. رابطه تجربی \eqref{eq:mobility-n} مجموعه پارامترهای
دیگری دارد. برای مقایسه کنترل‌شده اثر \lr{Velocity Saturation} در این جزوه،
\lr{Mobility} مؤثر را در دو محاسبه یکسان می‌گیریم. تفسیر کاهش \lr{Mobility}
در \eqref{eq:mobility-n} بر اساس صورت ۵۴۰ و وابستگی به بایاس همان مدل
انجام می‌شود \cite[مثال‌های \secnum{۲.۱} و \secnum{۲.۳}]{weste}.

\subsubsection*{مدل سرعت و دلیل ضریب دو در میدان بحرانی}
در میدان کوچک، سرعت رانش با میدان متناسب است و \(v=\mu_{\mathrm{eff}}E\).
در میدان بزرگ، سرعت متوسط به \(v_{\mathrm{sat}}\) نزدیک می‌شود. مدل قطعه‌ای
کتاب این رفتار را چنین تقریب می‌زند
\cite[روابط \secnum{۲.۲۴} و \secnum{۲.۲۵}]{weste}\cite[اسلاید ۹]{nonideal}
\begin{equation}
 v(E)=\begin{cases}
 \dfrac{\mu_{\mathrm{eff}}E}{1+E/E_c},&0\le E<E_c,\\[5pt]
 v_{\mathrm{sat}},&E\ge E_c.
 \end{cases}
 \label{eq:velocity-model}
\end{equation}
نماد \(E_c\) میدان بحرانی است. در مرز \(E=E_c\)، حد شاخه اول برابر
\(\mu_{\mathrm{eff}}E_c/2\) است. برای پیوستگی سرعت، این مقدار باید با
\(v_{\mathrm{sat}}\) برابر باشد. ازاین‌رو
\begin{equation}
 \frac{\mu_{\mathrm{eff}}E_c}{2}=v_{\mathrm{sat}}
 \quad\Longrightarrow\quad
 \boxed{E_c=\frac{2v_{\mathrm{sat}}}{\mu_{\mathrm{eff}}}}.
 \label{eq:critical-field}
\end{equation}
ضریب دو از شکل مشخص شاخه اول در \eqref{eq:velocity-model} به دست آمده است.
برای \(E\ll E_c\)، مخرج تقریباً یک است و قانون میدان کوچک بازیابی می‌شود.
نمودار بی‌بعد شکل \ref{fig:velocity-normalized} همین اتصال دو شاخه را نشان
می‌دهد. با تعریف \(u=E/E_c\)، شاخه اول برابر
\(v/v_{\mathrm{sat}}=2u/(1+u)\) می‌شود.

\begin{figure}[!htbp]
\centering
\begin{latin}
\PersianMathDigits
\begin{tikzpicture}[x=3.2cm,y=2.5cm]
 \draw[->] (0,0)--(2.15,0) node[right] {\(E/E_c\)};
 \draw[->] (0,0)--(0,2.15) node[above] {\(v/v_{\mathrm{sat}}\)};
 \draw[densely dashed,gray] (0,0)--(1,2) node[above right] {\(v=\mu_{\mathrm{eff}}E\)};
 \draw[thick,domain=0:1,samples=70] plot (\x,{2*\x/(1+\x)});
 \draw[thick] (1,1)--(2,1);
 \draw[dotted] (1,0)--(1,1)--(0,1);
 \foreach \x in {0,1,2} {\draw (\x,0)--(\x,-0.03) node[below] {\(\x\)};}
 \foreach \y in {1,2} {\draw (0,\y)--(-0.03,\y) node[left] {\(\y\)};}
\end{tikzpicture}
\end{latin}
\caption{نمودار آموزشی مدل سرعت \eqref{eq:velocity-model}. خط‌چین امتداد قانون میدان کوچک و خط پیوسته مدل قطعه‌ای کتاب را نشان می‌دهد.}
\label{fig:velocity-normalized}
\end{figure}

در کتاب، \(v_{\mathrm{sat},n}\approx\qty{e7}{\centi\metre\per\second}\)
و \(v_{\mathrm{sat},p}\approx\qty{8e6}{\centi\metre\per\second}\) ذکر
شده‌اند \cite[صفحه چاپی ۷۶]{weste}. این مقادیر و ضرایب \lr{Mobility}،
پارامترهای مدل مورد استفاده مرجع‌اند. در اسلاید ۹، استقلال \(E_c\) از
\(V_{DS}\) و \(L\) در مدل بیان شده است. با ثابت بودن بایاس گیت و پارامترهای
ماده، \eqref{eq:critical-field} نیز همین رفتار را دارد. هنگام تغییر
\(V_{GS}\)، \(\mu_{\mathrm{eff}}\) می‌تواند تغییر کند و مقدار \(E_c\)
باید سازگار با آن به‌روزرسانی شود.

\subsubsection*{ولتاژ بحرانی و حل مثال \secnum{۲.۴}}
ولتاژ بحرانی را با استفاده از طول کانال تعریف می‌کنیم
\begin{equation}
 V_c=E_cL=\frac{2v_{\mathrm{sat}}L}{\mu_{\mathrm{eff}}}.
 \label{eq:critical-voltage}
\end{equation}
واحد \(E_cL\) برابر ولت است. در مثال کتاب، \(L=\qty{50}{\nano\metre}
=\qty{5e-6}{\centi\metre}\) است. جایگذاری مقادیر گرد‌شده
\eqref{eq:mobility-example-results} می‌دهد
\begin{align}
 E_{c,n}&=\frac{2\times10^7}{96}
 \approx\qty{2.08e5}{\volt\per\centi\metre},
 &V_{c,n}&\approx(2.08\times10^5)(5\times10^{-6})\,\unit{\volt}
 \approx\qty{1.04}{\volt},\notag\\
 E_{c,p}&=\frac{2(8\times10^6)}{36}
 \approx\qty{4.44e5}{\volt\per\centi\metre},
 &V_{c,p}&\approx\qty{2.22}{\volt}.
 \label{eq:critical-example}
\end{align}
این نتایج با مثال \secnum{۲.۴} کتاب و اسلاید ۱۲ تطبیق دارند
\cite{weste,nonideal}. در این نمونه، \(V_c\) در \lr{pMOS} بزرگ‌تر است و
برای اضافه‌ولتاژ گیت یکسان، اثر محدودشدن سرعت ضعیف‌تر خواهد بود. نسبت
\(V_{GT}/V_c\) در جریان استخراج‌شده این تفاوت را کمّی می‌کند.

تعریف \(V_c\) با میدان متوسط، مقیاسی برای مقایسه بایاس‌هاست. میدان محلی
نزدیک \lr{Drain} می‌تواند زودتر از میدان متوسط به \(E_c\) برسد.
ولتاژ شروع اشباع جریان، \(V_{\mathrm{dsat}}\)، در ادامه از شرط مرزی همان
مدل تعیین می‌شود.

\subsubsection*{استخراج رابطه جریان با میدان محلی}
محور \(x\) را از \lr{Source} به \lr{Drain} می‌گیریم و \(V(x)\) را
پتانسیل کانال نسبت به \lr{Source} تعریف می‌کنیم. کران‌ها
\(V(0)=0\) و \(V(L)=V_{DS}\) هستند. برای اختصار،
\(A=V_{GT}>0\) و \(D=V_{DS}\) می‌گذاریم. بزرگی بار متحرک بر واحد
طول کانال در مدل تدریجی برابر است با
\begin{equation}
 q'(x)=WC_{\mathrm{ox}}[A-V(x)].
 \label{eq:local-charge}
\end{equation}
این بیان بار، همان مبنای بخش \secnum{۲.۲} کتاب است. تغییر این استخراج،
استفاده از رابطه سرعت \eqref{eq:velocity-model} است
\cite[بخش‌های \secnum{۲.۲} و \secnum{۲.۴.۱}]{weste}.

پیش از رسیدن میدان انتهای کانال به \(E_c\)، با
\(E(x)=\mathrm{d}V/\mathrm{d}x\)، جریان در هر مقطع چنین نوشته می‌شود
\begin{equation}
 I=WC_{\mathrm{ox}}(A-V)
 \frac{\mu_{\mathrm{eff}}\,\mathrm{d}V/\mathrm{d}x}
 {1+(\mathrm{d}V/\mathrm{d}x)/E_c}.
 \label{eq:current-local}
\end{equation}
در محاسبه برای یک بایاس گیت ثابت، \(\mu_{\mathrm{eff}}\) را در طول کانال
ثابت می‌گیریم. جریان \lr{DC} نیز در همه مقاطع یکسان است. ابتدا دو طرف را
در مخرج و سپس در \(\mathrm{d}x\) ضرب می‌کنیم
\begin{align}
 I+\frac{I}{E_c}\frac{\mathrm{d}V}{\mathrm{d}x}
 &=\mu_{\mathrm{eff}}WC_{\mathrm{ox}}(A-V)
 \frac{\mathrm{d}V}{\mathrm{d}x},\notag\\
 I\,\mathrm{d}x
 &=\left[\mu_{\mathrm{eff}}WC_{\mathrm{ox}}(A-V)-\frac{I}{E_c}\right]
 \mathrm{d}V.
 \label{eq:current-separated}
\end{align}
اکنون سمت چپ روی طول کانال و سمت راست روی اختلاف پتانسیل جمع می‌شود.
با استفاده از انتگرال‌های مرور پیش‌نیاز داریم
\begin{align}
 IL&=\mu_{\mathrm{eff}}WC_{\mathrm{ox}}
 \int_0^D(A-V)\,\mathrm{d}V-\frac{I}{E_c}\int_0^D\mathrm{d}V,\notag\\
 IL&=\mu_{\mathrm{eff}}WC_{\mathrm{ox}}\left(AD-\frac{D^2}{2}\right)
 -\frac{ID}{E_c},\notag\\
 I\left(L+\frac{D}{E_c}\right)
 &=\mu_{\mathrm{eff}}WC_{\mathrm{ox}}\left(AD-\frac{D^2}{2}\right).
 \label{eq:current-integrated}
\end{align}
با تقسیم بر \(L\)، استفاده از \(E_cL=V_c\) و تعریف
\(\beta_{\mathrm{eff}}=\mu_{\mathrm{eff}}C_{\mathrm{ox}}W/L\)، نتیجه می‌شود
\begin{equation}
 \boxed{I_{\mathrm{lin}}=
 \frac{\beta_{\mathrm{eff}}\left(V_{GT}V_{DS}-V_{DS}^2/2\right)}
 {1+V_{DS}/V_c}},\qquad 0\le V_{DS}<V_{\mathrm{dsat}}.
 \label{eq:high-field-linear}
\end{equation}
این رابطه با شاخه نخست رابطه \secnum{۲.۲۶} کتاب برابر است. مخرج بی‌بعد،
اثر کاهش سرعت نسبت به قانون میدان کوچک را در جریان وارد می‌کند. اگر
\(V_{DS}/V_c\) کوچک باشد، مخرج به یک نزدیک می‌شود.

\subsubsection*{شرط انتهای کانال و اثبات تمرین \secnum{۲.۶}}
تمرین \secnum{۲.۶}، اثبات رابطه \secnum{۲.۲۷} کتاب را می‌خواهد
\cite[صفحه چاپی ۹۶]{weste}. در مرز اشباع این مدل، میدان محلی انتهای کانال
به \(E_c\) و سرعت به \(v_{\mathrm{sat}}\) می‌رسد. اگر
\(s=V_{\mathrm{dsat}}\) باشد، رابطه \eqref{eq:local-charge} بار بر واحد
طول انتهای کانال را برابر \(WC_{\mathrm{ox}}(A-s)\) می‌دهد. بنابراین
\begin{equation}
 I_{\mathrm{sat}}=WC_{\mathrm{ox}}(A-s)v_{\mathrm{sat}}.
 \label{eq:drain-saturation-current}
\end{equation}
در این مدل، برای \(V_{DS}\ge s\)، افزایش بایاس \lr{Drain} با ناحیه
میدان قوی انتهای کانال همراه است و جریان را به همین مقدار محدود می‌کنیم
\cite[رابطه \secnum{۲.۲۶}]{weste}.

در مرز، دو شاخه باید یک جریان بدهند. از \eqref{eq:critical-voltage}،
\(v_{\mathrm{sat}}=\mu_{\mathrm{eff}}V_c/(2L)\) است. جایگذاری می‌دهد
\begin{equation}
 \frac{\mu_{\mathrm{eff}}WC_{\mathrm{ox}}}{L}
 \frac{(A-s/2)s}{1+s/V_c}
 =WC_{\mathrm{ox}}(A-s)\frac{\mu_{\mathrm{eff}}V_c}{2L}.
 \label{eq:boundary-equality}
\end{equation}
عامل \(\mu_{\mathrm{eff}}WC_{\mathrm{ox}}/L\) در دو طرف مثبت و مشترک است.
آن را حذف می‌کنیم و دو طرف را در \(2(1+s/V_c)\) ضرب می‌کنیم
\begin{align}
 2(A-s/2)s&=(A-s)(V_c+s),\notag\\
 2As-s^2&=AV_c+As-V_cs-s^2,\notag\\
 As+V_cs&=AV_c,\notag\\
 (A+V_c)s&=AV_c.
 \label{eq:saturation-algebra}
\end{align}
جمله‌های \(-s^2\) از دو طرف حذف می‌شوند. چون \(A+V_c>0\)، تقسیم بر آن
مجاز است و جواب یکتا به دست می‌آید
\begin{equation}
 \boxed{V_{\mathrm{dsat}}=\frac{V_{GT}V_c}{V_{GT}+V_c}}.
 \label{eq:high-field-saturation-voltage}
\end{equation}
صورت، واحد ولت به توان دو و مخرج، واحد ولت دارد. پس جواب ولتاژ است.
برای \(V_{GT}>0,V_c>0\)، دو نسبت زیر را می‌نویسیم
\begin{equation}
 \frac{V_{\mathrm{dsat}}}{V_{GT}}=\frac{V_c}{V_{GT}+V_c}<1,\qquad
 \frac{V_{\mathrm{dsat}}}{V_c}=\frac{V_{GT}}{V_{GT}+V_c}<1.
 \label{eq:saturation-bounds}
\end{equation}
در نتیجه، \(0<V_{\mathrm{dsat}}<\min(V_{GT},V_c)\) و بار انتهای کانال
در مرز مثبت است. شکل جبری \(ab/(a+b)\) در فرمول دو مقاومت موازی نیز
دیده می‌شود. در اینجا دو عامل \(V_{GT},V_c\) هر دو ولتاژند.

با جایگذاری \eqref{eq:high-field-saturation-voltage} در
\eqref{eq:drain-saturation-current}، ابتدا اختلاف \(A-s\) را ساده می‌کنیم
\[
 A-s=A-\frac{AV_c}{A+V_c}=\frac{A^2}{A+V_c}.
\]
بنابراین جریان اشباع به شکل زیر نوشته می‌شود
\begin{equation}
 \boxed{I_{\mathrm{sat}}
 =WC_{\mathrm{ox}}v_{\mathrm{sat}}\frac{V_{GT}^2}{V_{GT}+V_c}}.
 \label{eq:high-field-saturation-current}
\end{equation}
این نتیجه همان رابطه \secnum{۲.۲۸} کتاب است. ضرب
\(WC_{\mathrm{ox}}\) در یک ولتاژ و یک سرعت، واحد آمپر می‌دهد.

\subsubsection*{حدهای مدل و تفاوت دو معنای \lr{Saturation}}
برای آشکارشدن حد کانال بلند، از \(v_{\mathrm{sat}}=
\mu_{\mathrm{eff}}V_c/(2L)\) استفاده می‌کنیم
\begin{equation}
 I_{\mathrm{sat}}=
 \frac{\beta_{\mathrm{eff}}V_{GT}^2}{2}
 \frac{1}{1+V_{GT}/V_c}.
 \label{eq:current-reduction-factor}
\end{equation}
اگر \(V_{GT}\ll V_c\)، عامل آخر به یک و
\(V_{\mathrm{dsat}}\) به \(V_{GT}\) نزدیک می‌شود. جریان مربعی مدل
\lr{Long-channel} با همان \lr{Mobility} مؤثر به دست می‌آید. اگر
\(V_{GT}\gg V_c\)، با تقسیم صورت و مخرج
\eqref{eq:high-field-saturation-current} بر \(V_{GT}\) داریم
\begin{equation}
 I_{\mathrm{sat}}=
 WC_{\mathrm{ox}}v_{\mathrm{sat}}V_{GT}
 \frac{1}{1+V_c/V_{GT}}
 \approx WC_{\mathrm{ox}}v_{\mathrm{sat}}V_{GT}.
 \label{eq:velocity-saturated-limit}
\end{equation}
در این حد، جریان تقریباً با \(V_{GT}\) متناسب و
\(V_{\mathrm{dsat}}\approx V_c\) است
\cite[روابط \secnum{۲.۲۸} و \secnum{۲.۲۹}]{weste}.

واژه \lr{Saturation} برای ناحیه کاری جریان به کار می‌رود.
\lr{Velocity Saturation} محدودشدن سرعت حامل‌ها را توصیف می‌کند.
مدل کانال بلند نیز ناحیه \lr{Saturation} دارد. در مدل حاضر، میدان بحرانی
محلی شرط مرز را تعیین کرده است و می‌تواند جریان را در ولتاژی کوچک‌تر از
\(V_{GT}\) به اشباع برساند \cite[پانوشت صفحه چاپی ۷۵]{weste}.

برای استفاده از \eqref{eq:high-field-linear}، ابتدا باید
\(V_{\mathrm{dsat}}\) حساب شود. شاخه \lr{Linear} فقط تا این مرز معتبر
است. اتصال دو شاخه مدل قطعه‌ای پیوستگی جریان را برقرار می‌کند. شیب دو
شاخه در مرز می‌تواند متفاوت باشد، زیرا مدل سرعت نیز در \(E_c\) اتصال
قطعه‌ای دارد. تعیین مرز از شرط میدان انتهای کانال با همین فرض‌ها سازگار است.

\subsubsection*{مثال آموزشی جریان و کنترل ناحیه کاری}
پارامترهای هندسی و اکسید مثال \secnum{۲.۱} کتاب را با \lr{Mobility} مؤثر
گرد‌شده مثال \secnum{۲.۳} به کار می‌بریم. فرض‌ها در جدول
\ref{tab:high-field-data} آمده‌اند. این محاسبه، کاربرد آموزشی مدل
\eqref{eq:high-field-linear} و \eqref{eq:high-field-saturation-current} است.

\begin{table}[!htbp]
\centering
\caption{پارامترهای مثال آموزشی جریان. بایاس بدنه و آستانه ثابت فرض شده‌اند.}
\label{tab:high-field-data}
\begin{tabular}{cc@{\qquad}cc}
\toprule
کمیت & مقدار & کمیت & مقدار\\
\midrule
\(W\) & \(\qty{100}{\nano\metre}\) & \(L\) & \(\qty{50}{\nano\metre}\)\\
\(t_{\mathrm{ox}}\) & \(\qty{1.05}{\nano\metre}\) & \(V_t\) & \(\qty{0.3}{\volt}\)\\
\(V_{GS}\) & \(\qty{1.0}{\volt}\) & \(v_{\mathrm{sat}}\) & \(\qty{e5}{\metre\per\second}\)\\
\(\mu_{\mathrm{eff}}\) & \(\qty{0.0096}{\metre\squared\per\volt\per\second}\) & \(\varepsilon_{\mathrm{ox}}\) & \(3.9\varepsilon_0\)\\
\bottomrule
\end{tabular}
\end{table}

با ثابت \lr{Permittivity} خلأ
\(\varepsilon_0\approx\qty{8.85e-12}{\farad\per\metre}\)، نخست
ظرفیت بر واحد سطح و ضریب جریان را می‌یابیم
\begin{align}
 C_{\mathrm{ox}}&=\frac{3.9(8.85\times10^{-12})}{1.05\times10^{-9}}
 \approx\qty{0.03287}{\farad\per\metre\squared},\notag\\
 \beta_{\mathrm{eff}}&=(0.0096)(0.03287)\frac{100}{50}
 \approx\qty{631.1}{\micro\ampere\per\volt\squared},\notag\\
 V_{GT}&=\qty{0.7}{\volt},\qquad
 V_c\approx\qty{1.0417}{\volt},\qquad
 V_{\mathrm{dsat}}=\frac{0.7(1.0417)}{0.7+1.0417}\,\unit{\volt}
 \approx\qty{0.4187}{\volt}.
 \label{eq:high-field-example-parameters}
\end{align}
برای \(V_{DS}=\qty{0.1}{\volt}\)، شرط شاخه \lr{Linear} برقرار است.
جریان و عامل کاهش آن برابرند با
\begin{equation}
 I_{DS}\approx\frac{631.1[0.7(0.1)-(0.1)^2/2]}{1+0.1/1.0417}
 \,\unit{\micro\ampere}
 \approx\qty{37.4}{\micro\ampere}.
 \label{eq:high-field-example-linear}
\end{equation}
اگر در همین بایاس مخرج را یک بگیریم، جریان حدود
\(\qty{41.0}{\micro\ampere}\) می‌شود. نسبت دو جریان
\(1/(1+0.1/1.0417)\approx0.912\) است.

برای \(V_{DS}=\qty{1.0}{\volt}\)، شرط
\(V_{DS}>V_{\mathrm{dsat}}\) برقرار است. از رابطه اشباع به دست می‌آوریم
\begin{equation}
 I_{DS}\approx(100\times10^{-9})(0.03287)(10^5)
 \frac{0.7^2}{0.7+1.0417}\,\unit{\ampere}
 \approx\qty{92.5}{\micro\ampere}.
 \label{eq:high-field-example-sat}
\end{equation}
مرجع مقایسه با همان \(\mu_{\mathrm{eff}}\)، جریان مربعی
\(\beta_{\mathrm{eff}}V_{GT}^2/2\approx\qty{154.6}{\micro\ampere}\)
است. نسبت \eqref{eq:current-reduction-factor} حدود \(0.598\) می‌شود.
پس محدودشدن سرعت در این مقایسه کنترل‌شده، جریان را حدود \num{40.2} درصد
کاهش می‌دهد. تغییرات \lr{Mobility} و اثرهای آستانه در این مقایسه ثابت‌اند.

شکل \ref{fig:high-field-current} دو شاخه مدل همین مثال را نشان می‌دهد.
افزایش ولتاژ در ناحیه اشباع این مدل، جریان را ثابت نگه می‌دارد. در بخش
\ref{sec:clm}، اثر افزایش طول ناحیه تخلیه بر شیب جریان بررسی می‌شود.
\begin{figure}[!htbp]
\centering
\begin{latin}
\PersianMathDigits
\begin{tikzpicture}[x=8cm,y=0.035cm]
 \draw[->] (0,0)--(1.1,0) node[right] {\(V_{DS}\;(\mathrm{V})\)};
 \draw[->] (0,0)--(0,175) node[above] {\(I_{DS}\;(\mathrm{\mu A})\)};
 \draw[densely dashed,gray,domain=0:0.7,samples=80] plot (\x,{631.1314*(0.7*\x-\x*\x/2)});
 \draw[densely dashed,gray] (0.7,154.6272)--(1,154.6272);
 \draw[thick,domain=0:0.4186603,samples=80] plot (\x,{631.1314*(0.7*\x-\x*\x/2)/(1+\x/1.0416667)});
 \draw[thick] (0.4186603,92.4804)--(1,92.4804);
 \draw[dotted] (0.4186603,0)--(0.4186603,92.4804);
 \node[above] at (0.82,154.6272) {\(\mu_{\mathrm{eff}}\) fixed};
 \node[above] at (0.78,92.4804) {velocity model};
 \foreach \x in {0,0.5,1} {\draw (\x,0)--(\x,-2) node[below] {\(\x\)};}
 \foreach \y in {50,100,150} {\draw (0,\y)--(-0.012,\y) node[left] {\(\y\)};}
\end{tikzpicture}
\end{latin}
\caption{مقایسه آموزشی مدل میدان قوی با مدل مربعی، با \lr{Mobility} مؤثر یکسان و پارامترهای جدول \ref{tab:high-field-data}.}
\label{fig:high-field-current}
\end{figure}

\subsubsection*{مدل \lr{Alpha-power law} و معنای پارامترهای آن}
برای برازش ساده‌تر منحنی‌ها، کتاب و اسلاید ۱۱ مدل \lr{Alpha-power law}
را معرفی می‌کنند. \(\alpha\) توان تجربی وابستگی جریان به اضافه‌ولتاژ گیت
است. در توصیف مرجع، مقدار آن از حدود دو در رفتار کانال بلند به حدود یک
در محدودیت شدید سرعت تغییر می‌کند
\cite[روابط \secnum{۲.۳۰} و \secnum{۲.۳۱}]{weste}\cite[اسلایدهای ۱۰ و ۱۱]{nonideal}.
با \(V_{GT}>0\)، دو پارامتر شاخه‌ها چنین تعریف می‌شوند
\begin{equation}
 I_{\mathrm{dsat}}=\frac{\beta P_c}{2}V_{GT}^{\alpha},\qquad
 V_{\mathrm{dsat},\alpha}=P_vV_{GT}^{\alpha/2}.
 \label{eq:alpha-parameters}
\end{equation}
نماد \(P_c\) ضریب برازش جریان و \(P_v\) ضریب برازش ولتاژ است. جریان
کل با خط مستقیم در شاخه \lr{Linear} تقریب زده می‌شود
\begin{equation}
 I_{DS}=\begin{cases}
 0,&V_{GT}\le0,\\[2pt]
 I_{\mathrm{dsat}}\dfrac{V_{DS}}{V_{\mathrm{dsat},\alpha}},
 &V_{GT}>0,\quad 0\le V_{DS}<V_{\mathrm{dsat},\alpha},\\[5pt]
 I_{\mathrm{dsat}},&V_{GT}>0,\quad V_{DS}\ge V_{\mathrm{dsat},\alpha}.
 \end{cases}
 \label{eq:alpha-piecewise}
\end{equation}
در مرز دو شاخه، عامل \(V_{DS}/V_{\mathrm{dsat},\alpha}=1\) و جریان
پیوسته است. مقدار \(V_{\mathrm{dsat},\alpha}\) باید از پارامترهای همین
برازش محاسبه شود. مرز \eqref{eq:high-field-saturation-voltage} متعلق به
مدل سرعت بخش قبل است. کاربرد هر دو مرز در یک محاسبه به بازبرازش سازگار
نیاز دارد.

اگر \(\beta\) واحد \(\unit{\ampere\per\volt\squared}\) داشته باشد،
سازگاری ابعادی \eqref{eq:alpha-parameters} ایجاب می‌کند
\begin{equation}
 [P_c]=\mathrm{V}^{\,2-\alpha},\qquad
 [P_v]=\mathrm{V}^{\,1-\alpha/2}.
 \label{eq:alpha-dimensions}
\end{equation}
برای \(\alpha=2\)، هر دو ضریب بی‌بعد می‌شوند. برای توان دیگر، مقدار عددی
آن‌ها به قرارداد واحد ولتاژ وابسته است. می‌توان با یک ولتاژ مرجع مثبت
\(V_0\)، بیان بی‌بعد زیر را به کار برد
\begin{equation}
 I_{\mathrm{dsat}}=K_I\left(\frac{V_{GT}}{V_0}\right)^\alpha,\qquad
 V_{\mathrm{dsat},\alpha}=K_V
 \left(\frac{V_{GT}}{V_0}\right)^{\alpha/2}.
 \label{eq:alpha-normalized}
\end{equation}
در این بیان، \(K_I\) جریان و \(K_V\) ولتاژ هستند و
\(K_I=\beta P_cV_0^\alpha/2\)، \(K_V=P_vV_0^{\alpha/2}\).
این بازنویسی آموزشی، همان مدل را با واحدهای آشکار بیان می‌کند.

\subsubsection*{مثال آموزشی برازش توان و استفاده از مدل}
فرض کنید دو مقدار جریان اشباع برای یک افزاره فرضی در یک دما و هندسه
ثابت در اختیار داریم. در \(V_{GT,1}=\qty{0.4}{\volt}\)، جریان
\(I_1=\qty{120}{\micro\ampere}\) و در
\(V_{GT,2}=\qty{0.8}{\volt}\)، جریان
\(I_2=\qty{300}{\micro\ampere}\) است. فرض می‌کنیم هر دو نقطه در ناحیه
اشباع و در دامنه معتبر یک برازش مشترک قرار دارند. با تقسیم دو رابطه
جریان در \eqref{eq:alpha-normalized}، \(K_I\) حذف می‌شود
\begin{align}
 \frac{I_2}{I_1}&=\left(\frac{V_{GT,2}}{V_{GT,1}}\right)^\alpha,\notag\\
 \alpha&=\frac{\ln(I_2/I_1)}{\ln(V_{GT,2}/V_{GT,1})}
 =\frac{\ln(2.5)}{\ln2}\approx1.322.
 \label{eq:alpha-fit}
\end{align}
نسبت جریان‌ها و نسبت ولتاژها هر دو بی‌بعدند. برای \(V_0=\qty{1}{\volt}\)،
ضریب جریان از نقطه نخست به دست می‌آید.
\[
 K_I=\frac{120}{(0.4)^{1.322}}\approx\qty{402.9}{\micro\ampere}.
\]
پیش‌بینی همین برازش در \(V_{GT}=\qty{0.6}{\volt}\)،
حدود \(\qty{205.1}{\micro\ampere}\) است. برای تعیین مرز اشباع، هنوز
داده‌ای برای \(K_V\) لازم است.

در مثال فرضی مستقل دیگری، \(\alpha=1.3\)،
\(K_I=\qty{400}{\micro\ampere}\)، \(K_V=\qty{0.4}{\volt}\) و
\(V_0=\qty{1}{\volt}\) را معلوم می‌گیریم. در
\(V_{GT}=\qty{0.5}{\volt}\)، داریم
\begin{align}
 I_{\mathrm{dsat}}&=400(0.5)^{1.3}\,\unit{\micro\ampere}
 \approx\qty{162.5}{\micro\ampere},\notag\\
 V_{\mathrm{dsat},\alpha}&=0.4(0.5)^{0.65}\,\unit{\volt}
 \approx\qty{0.2549}{\volt}.
 \label{eq:alpha-example}
\end{align}
اگر \(V_{DS}=\qty{0.1}{\volt}\) باشد، جریان شاخه \lr{Linear} برابر
\(162.5(0.1/0.2549)\approx\qty{63.7}{\micro\ampere}\) است.
در \(V_{DS}=\qty{0.7}{\volt}\)، جریان برابر مقدار اشباع
\(\qty{162.5}{\micro\ampere}\) می‌شود. شکل \secnum{۲.۱۶} کتاب نشان
می‌دهد که دقت این تقریب خطی در ولتاژهای کوچک \lr{Drain} محدود است
\cite[صفحه چاپی ۷۸]{weste}. ضرایب این دو مثال آموزشی از داده یک فناوری
واقعی استخراج نشده‌اند.

\subsubsection*{جریان روشن، ولتاژ تغذیه و تقریب خطی ناحیه روشن}
جریان روشن در اسلاید ۵ با بایاس \(V_{GS}=V_{DS}=V_{DD}\) تعریف شده
است. در مدل توان، با ثابت بودن ضرایب و آستانه و برقرار بودن اشباع، داریم
\begin{equation}
 I_{\mathrm{on}}\propto(V_{DD}-V_t)^\alpha.
 \label{eq:on-current-alpha}
\end{equation}
عبارت \(I_{DS}\propto V_{DD}^{\alpha}\) در اسلاید ۱۰، وابستگی ساده‌شده
به تغذیه را بیان می‌کند. برای محاسبه عددی با آستانه معلوم، اضافه‌ولتاژ
\(V_{DD}-V_t\) در \eqref{eq:on-current-alpha} وارد می‌شود.
برای مثال آموزشی، کاهش تغذیه از \(\qty{1.0}{\volt}\) به
\(\qty{0.8}{\volt}\)، با \(V_t=\qty{0.3}{\volt}\) و \(\alpha=1.3\)،
نسبت جریان \((0.5/0.7)^{1.3}\approx0.646\) می‌دهد. نسبت
\((0.8/1.0)^{1.3}\approx0.748\) متناظر با حذف آستانه از این وابستگی است.

برای ارتباط با تأخیر گیت، ظرفیت بار خروجی را \(C_L\) و جریان متوسط
گذار را \(I_{\mathrm{avg}}\) می‌نامیم. برآورد
\(t_{\mathrm{sw}}\approx C_LV_{DD}/I_{\mathrm{avg}}\) از تقسیم بار
\(C_LV_{DD}\) بر جریان متوسط به دست می‌آید
\cite[بخش \secnum{۱.۹}]{weste}. اگر ظرفیت بار ثابت باشد و نسبت جریان
روشن را تقریب نسبت جریان متوسط گذار بگیریم، در همین مثال نسبت تأخیر
\(0.8/0.646\approx1.24\) می‌شود. کاهش دامنه ولتاژ و کاهش جریان در
این برآورد با یکدیگر رقابت می‌کنند و تأخیر حدود ۲۴ درصد افزایش می‌یابد.
این محاسبه به فرض ثابت‌بودن بار، آستانه و ضرایب برازش و تقریب جریان
متوسط وابسته است.

کتاب برای ناحیه به‌اندازه کافی روشن، یک تقریب خطی دیگر نیز معرفی می‌کند
\cite[رابطه \secnum{۲.۳۲} و شکل \secnum{۲.۱۷}]{weste}
\begin{equation}
 I_{DS}\approx k(V_{GS}-V_t^*),\qquad V_{DS}=V_{GS}.
 \label{eq:linear-on-fit}
\end{equation}
در این رابطه، \(k\) شیب برازش با واحد
\(\unit{\ampere\per\volt}\) و \(V_t^*\) محل برخورد امتداد خط با محور
ولتاژ است. این کمیت از برازش ناحیه جریان زیاد تعیین می‌شود.
آستانه \(V_t\) در مدل تشکیل کانال جداگانه تعریف شده است.

برای مثال آموزشی، فرض کنید جریان در \(V_{GS}=\qty{0.7}{\volt}\) برابر
\(\qty{0.48}{\milli\ampere}\) و در
\(V_{GS}=\qty{0.9}{\volt}\) برابر \(\qty{0.72}{\milli\ampere}\) باشد.
در هر دو نقطه، \(V_{DS}=V_{GS}\) است. شیب و عرض از محور ولتاژ چنین‌اند
\begin{equation}
 k=\frac{0.72-0.48}{0.9-0.7}
 =\qty{1.2}{\milli\ampere\per\volt},\qquad
 V_t^*=0.7-\frac{0.48}{1.2}=\qty{0.3}{\volt}.
 \label{eq:linear-on-example}
\end{equation}
پیش‌بینی خط در \(V_{GS}=\qty{0.8}{\volt}\)، جریان
\(\qty{0.60}{\milli\ampere}\) است. اعتبار این برازش به همان بازه بایاس
محدود است. در نزدیکی خاموشی، مدل جریان زیرآستانه باید جداگانه بررسی شود.
